% Per-volume bibliography for the CAS PCBA 8 monograph.
% Only works actually cited in this volume appear here; the program-wide
% shared list (../shared/bibliography.tex) serves the sibling monographs.
\begin{thebibliography}{99}
\bibitem{harel} D.~Harel, ``Statecharts: a visual formalism for complex systems,'' \emph{Sci.\ Comput.\ Program.}\ 8 (1987), 231--274.
\bibitem{vhdl} IEEE Std 1076-2019, \emph{VHDL Language Reference Manual}.
\bibitem{ieee754} IEEE Std 754-2019, \emph{Floating-Point Arithmetic}.
\bibitem{tarjan} R.~Tarjan, ``Depth-first search and linear graph algorithms,'' \emph{SIAM J.\ Comput.}\ 1 (1972), 146--160.
\bibitem{quickcheck} K.~Claessen and J.~Hughes, ``QuickCheck: a lightweight tool for random testing of Haskell programs,'' \emph{ICFP} (2000).
\bibitem{claessen} K.~Claessen, ``Shrinking and showing functions,'' \emph{Haskell Symposium}, 2012.
\bibitem{fuzzchick} L.~Lampropoulos, M.~Hicks, and B.~C. Pierce, ``Coverage guided, property based testing,'' \emph{Proc.\ ACM Program.\ Lang.}\ 3 (OOPSLA), 2019.
\bibitem{hypothesis} D.~R. MacIver, Z.~Hatfield-Dodds, et al., ``Hypothesis: a new approach to property-based testing,'' \emph{J.\ Open Source Softw.}\ 4(43), 2019.
\bibitem{sargent} R.~G. Sargent, ``Verification and validation of simulation models,'' \emph{J.\ Simulation} 7 (2013), 12--24.
\bibitem{pearl} J.~Pearl, \emph{Causality: Models, Reasoning, and Inference}, 2nd ed., Cambridge University Press, 2009.
\bibitem{prov} W3C, \emph{PROV-O: The PROV Ontology}, 2013.
\bibitem{fisher} R.~A. Fisher, ``On the mathematical foundations of theoretical statistics,'' \emph{Philos.\ Trans.\ R.\ Soc.\ A} 222 (1922), 309--368.
\bibitem{rao} C.~R. Rao, ``Information and the accuracy attainable in the estimation of statistical parameters,'' \emph{Bull.\ Calcutta Math.\ Soc.}\ 37 (1945), 81--91.
\bibitem{amari} S.-i.~Amari, \emph{Information Geometry and Its Applications}, Springer, 2016.
\bibitem{ay} N.~Ay, J.~Jost, H.~V.~L\^e, and L.~Schwachh\"ofer, \emph{Information Geometry}, Springer, 2017.
\bibitem{merkle} R.~C. Merkle, ``A digital signature based on a conventional encryption function,'' \emph{CRYPTO}, 1987.
\bibitem{mco} Mars Climate Orbiter Mishap Investigation Board, \emph{Phase I Report}, NASA, 1999.
\end{thebibliography}
